\documentclass[UTF8,11pt]{ctexart}
\usepackage[a4paper,margin=25mm]{geometry}
\usepackage{amsmath,amssymb,amsthm,mathtools}
\usepackage[all]{xy}
\usepackage{xurl}
\usepackage{hyperref}
\usepackage{enumitem}
\hypersetup{hidelinks,breaklinks=true}
\setlength{\emergencystretch}{3em}
\newtheorem*{theorem}{Theorem}
\newtheorem*{lemma}{Lemma}
\newtheorem*{proposition}{Proposition}
\newtheorem*{corollary}{Corollary}
\newtheorem*{definition}{Definition}
\newcommand{\Spec}{\operatorname{Spec}}
\newcommand{\Hom}{\operatorname{Hom}}
\newcommand{\End}{\operatorname{End}}
\newcommand{\Gal}{\operatorname{Gal}}
\newcommand{\Cl}{\operatorname{Cl}}
\newcommand{\vol}{\operatorname{vol}}
\newcommand{\covol}{\operatorname{covol}}
\newcommand{\tr}{\operatorname{tr}}
\newcommand{\im}{\operatorname{im}}
\newcommand{\id}{\operatorname{id}}
\newcommand{\coker}{\operatorname{coker}}
\newcommand{\stacksref}[2]{\href{https://stacks.math.columbia.edu/tag/#1}{#2 [#1]}}
\begin{document}
\section*{007\quad 有限域Kakeya集的体量下界}
\subsection*{来源与整理范围}
集体来源：MIT OpenCourseWare，MIT 18.S997 \emph{The Polynomial Method}，Fall 2012；授课教师 Lawrence D. Guth。课程页面注明讲义由 MIT 学生提供并获准使用，所取文件未单列执笔人，故不将转录正文擅自署为 Guth 独著。本文收录的是该课程的人类讲义版本，不是研究论文原版。获取日期：2026-09-28。\par 原文：Lecture 3, \emph{The Finite-Field Nikodym and Kakeya Problems}, Theorem 0.2 及其完整证明，pp.~1--2；维数计数引用 Lecture 1, Corollaries 0.2--0.3，pp.~1--2。
\par\url{https://ocw.mit.edu/courses/18-s997-the-polynomial-method-fall-2012/1da7feb132bb4a859fd021f5cfec98cc_MIT18_S997F12_lec3.pdf}
\par\url{https://ocw.mit.edu/courses/18-s997-the-polynomial-method-fall-2012/43c5d5700347be070289fe2be4cf30c7_MIT18_S997F12_lec1.pdf}
\par\url{https://ocw.mit.edu/courses/18-s997-the-polynomial-method-fall-2012/pages/lecture-notes/}

\subsection*{作者原命题与人类证明}
Let $F$ be a finite field with $q$ elements.

A set $K\subset F^n$ is called a Kakeya set if it contains a line in every direction. In other words, for every vector $a\in F^n\setminus\{0\}$, there is a vector $b$ so that the line $\{at+b\mid t\in F\}$ is contained in $K$. A trivial example of a Kakeya set is the entire vector space $F^n$. Can one find a Kakeya set significantly smaller than this?

\begin{theorem}[0.2]
A Kakeya set $K\subset F^n$ has at least $c_nq^n$ elements.
\end{theorem}
\begin{proof}
Let $K$ be a Kakeya set. If $K$ is smaller than the conclusion, let $P$ be a polynomial vanishing on $K$ of degree $<q$. Let $d$ be the degree of $P$. Write $P=P_d+Q$, where $P_d$ is the sum of monomials of degree $d$ and $Q$ is a polynomial of degree $\leq d-1$.

Let $a$ be any non-zero vector. Choose $b$ so that the line $\{at+b\mid t\in F\}$ is contained in $K$. Consider the polynomial in one variable $R(t):=P(at+b)$. The polynomial $R$ vanishes for each $t\in F$. It has degree $\leq d<q$, and so its coefficeints all vanish. In particular, its coefficient of degree $d$ vanishes. But the coefficient of $t^d$ in $R$ is exactly $P_d(a)$. So we see that $P_d(a)$ vanishes for all $a\in F^n\setminus\{0\}$. But since the degree of $P_d$ is $d<q$, it easily follows that $P_d$ vanishes at $0$ also. Then we see that $P_d$ is identically zero, and we reach a contradiction.
\end{proof}


\subsection*{整理说明（非作者正文）}
本条采用作者原结论：存在只依赖维数的正常数 $c_n$，使 $|K|\geq c_nq^n$。\textbf{不是}种子题要求的 DKSS 界 $\bigl(q/(2-1/q)\bigr)^n$，不声称本证明覆盖该更强常数。

原证明中的消失多项式来自所引 Lecture 1 的非零多项式维数计数；所用另一标准先修事实为：总次数 $<q$ 的多项式若在 $F^n$ 的所有点取零，则为零多项式（Lecture 3，p.~1 的未编号引理）。两项均明确列作先修，不以该 Kakeya 定理本身作先修。下列维数计数原文供定位；不另计为题目。$V(d)$ 为 $n$ 元总次数至多 $d$ 的多项式向量空间，$s$ 为给定点数。

拼写 ``coefficeints'' 按原文保留；集合分隔符、空格与显示公式作 TeX 排版规范化。网页取得了 PDF 文本；该文件的页面截图接口返回错误，未声称逐页图像核对，也没有把转录稿冒充下载原件。

\paragraph{所引用的维数计数结果（Lecture 1，pp.~1--2）}
\begin{corollary}[0.2]
If $\dim V(d)>s$, then there is a non-zero polynomial $P$ of degree $\leq d$ that vanishes on the given finite set.
\end{corollary}
The dimension of $V(d)$ is $\binom{d+n}{n}\geq d^n/n!$. For $n$ fixed and $d$ large, $d^n/n!$ is a good approximation of the dimension. Therefore, we get the following corollary.
\begin{corollary}[0.3]
For any set of $s$ points in $F^n$, there is a non-zero polynomial that vanishes on the set with degree $\leq ns^{1/n}$.
\end{corollary}


\subsection*{可修改的简短标注（非作者正文）}
$c$：有限域点集的极值大小。

$a$：消失多项式的维数计数构造；直线限制与一元多项式根数界；最高次齐次部分；有限域多项式恒等判别。

\texttt{cross\_theory\_status = yes}。点集的方向覆盖条件通过直线限制转成最高次齐次多项式在每个方向上消失，再以次数约束导出点数下界。位置：Theorem 0.2 的证明，pp.~1--2。全部标注为非穷尽、可修改初稿。
\subsection*{许可与修改记录}
材料来自 MIT OpenCourseWare；原作者与课程见上文。按该站所示 \href{https://creativecommons.org/licenses/by-nc-sa/4.0/}{CC BY-NC-SA 4.0} 许可复用。本转录及其附加整理采用同一许可；不得据此声称作者或 MIT 认可本次整理。2026-09-28：ChatGPT 转录为 TeX，加入题号、来源、整理说明和初稿标注；数学正文保持作者语言及顺序。

\end{document}
